% Part: normal-modal-logic
% Chapter: tableaux
% Section: completeness

\documentclass[../../../include/open-logic-section]{subfiles}

\begin{document}

\olfileid{nml}{tab}{cpl}

\olsection{\Log{K}-এর সম্পূর্ণতা}

\begin{explain}
  !!{tableau}s পদ্ধতির সম্পূর্ণতা দেখাতে প্রমাণ করতে হবে:
  $\Gamma \Proves !A$ প্রতিষ্ঠার জন্য কোনো বন্ধ !!{tableau} না
  থাকলে $\Gamma \Entails/ !A$, অর্থাৎ একটি প্রতিমডেল আছে।
  কিন্তু ``কোনো বন্ধ !!{tableau} নেই'' মানে, তা নির্মাণের
  প্রত্যেক প্রচেষ্টাতেই অন্তত একটি শাখা খোলা থাকবে।
  মূল কৌশল হলো: একটি নির্দিষ্ট, \emph{পদ্ধতিগত ও পূর্ণাঙ্গ}
  নির্মাণেও এমন শাখা খোলা থাকবে; বন্ধ !!{tableau} থাকলে এই
  নির্মাণই তা পেত। এর খোলা শাখায় প্রতিমডেল নির্ধারণের
  সব তথ্য থাকবে। শাখায় ব্যবহৃত সব পূর্বসূচিই প্রতিমডেলের
  জগৎ হবে; আর !!a{propositional
    variable}~$p$,
  $\sigma$-তে সত্য হবে যদি এবং কেবল যদি
  $\sFmla{\True}{p}[\sigma]$ শাখায় থাকে।
\end{explain}

\begin{defn}
  !!a{tableau}-র একটি শাখাকে সম্পূর্ণ বলি, যদি প্রযোজ্য বিধির
  পূর্বসূচিযুক্ত !!{formula} $\sFmla{S}{!A}[\sigma]$
  শাখায় থাকলে
  নিম্নলিখিত সূত্রগুলিও শাখায় থাকে:
  \begin{enumerate}
    \item প্রস্তাবনামূলক স্তূপীকরণ-বিধির ক্ষেত্রে সংশ্লিষ্ট
      উপসংহারের সব পূর্বসূচিযুক্ত !!{formula}s;
    \item প্রস্তাবনামূলক শাখা-বিভাজন বিধির ক্ষেত্রে সংশ্লিষ্ট
      উপসংহারগুলির অন্তত একটি !!{formula}s;
    \item নতুন পূর্বসূচি-প্রয়োজনীয় মোডাল বিধির ক্ষেত্রে অন্তত
      একটি সম্ভাব্য উপসংহার;
    \item ব্যবহৃত পূর্বসূচি-প্রয়োজনীয় মোডাল বিধির ক্ষেত্রে
      শাখায় উপস্থিত প্রত্যেক প্রাসঙ্গিক পূর্বসূচির জন্য
      সংশ্লিষ্ট উপসংহার।
  \end{enumerate}
\end{defn}

\begin{explain}
যেমন, সম্পূর্ণ শাখায় $\sFmla{\True}{!B \land !C}[\sigma]$
থাকলে $\sFmla{\True}{!B}[\sigma]$ ও
$\sFmla{\True}{!C}[\sigma]$-ও থাকে।
% BN-SRC-377: উৎসের True-সংযোগ পূর্বপক্ষে পূর্বসূচি অনুপস্থিত।
$\sFmla{\True}{!B \lor !C}[\sigma]$ থাকলে
$\sFmla{\True}{!B}[\sigma]$ অথবা
$\sFmla{\True}{!C}[\sigma]$ অন্তত একটি থাকে।
% BN-SRC-378: True-বিয়োগের বাম উপসূত্রে উৎসের False চিহ্ন ভুল।
আর \iftag{prvBox}
{$\sFmla{\False}{\Box !B}[\sigma]$}
{$\sFmla{\True}{\Diamond !B}[\sigma]$}
থাকলে অন্তত একটি~$n$-এর জন্য \iftag{prvBox}
{$\sFmla{\False}{!B}[\sigma.n]$}
{$\sFmla{\True}{!B}[\sigma.n]$}-ও থাকে।
% BN-SRC-379: উৎসে মোডাল অন্তর্বর্তী সূত্র অনুপস্থিত; উত্তরসূরিতে অন্তর্বর্তী সূত্রই যায়।
অন্যদিকে \iftag{prvBox}
{$\sFmla{\True}{\Box !B}[\sigma]$}
{$\sFmla{\False}{\Diamond !B}[\sigma]$}
থাকলে শাখায় ব্যবহৃত প্রত্যেক $\sigma.n$-এর জন্য,
অর্থাৎ প্রাসঙ্গিক প্রত্যেক~$n$-এর জন্য,
\iftag{prvBox}{$\sFmla{\True}{!B}[\sigma.n]$}
{$\sFmla{\False}{!B}[\sigma.n]$}-ও থাকে।
% BN-SRC-380: ব্যবহৃত উত্তরসূরিতে একই চিহ্নের অন্তর্বর্তী সূত্র যায়।
\end{explain}

\begin{prop}\ollabel{prop:complete-tableau}
  প্রত্যেক সসীম $\Gamma$-র এমন একটি !!a{tableau} আছে যার
  প্রত্যেক খোলা শাখা সম্পূর্ণ।
  % BN-SRC-734: বদ্ধ শাখায় আর বিধি প্রয়োগের দাবি নেই।
\end{prop}

\begin{proof}
  $\Gamma$-র !!a{tableau}-র একটি খোলা শাখা নিই।
  ওই পর্যায়ে শাখার সসীমসংখ্যক পূর্বসূচিযুক্ত !!{formula}s-এ
  বিধি প্রয়োগ সম্ভব। একটি স্থির ক্রমে, ধরা যাক উপর থেকে
  নিচে, যেসব !!{formula}s-এ (1)--(4) শর্ত এখনও পূরণ হয়নি সেগুলিতে
  প্রযোজ্য বিধি প্রয়োগ করে শাখা বাড়াই। কোনো বিধিতে
  শাখা বিভক্ত হলে উৎপন্ন প্রতিটি শাখায় অবশিষ্ট !!{formula}s-এর
  জন্য একই কাজ করি। বর্তমান শাখার !!{formula}s ও ব্যবহৃত পূর্বসূচি
  উভয়ের সংখ্যাই সসীম, তাই এই পর্যায়ের কাজ শেষ হয়।
  তারপর উৎপন্ন শাখাগুলিতে একই প্রক্রিয়া পুনরাবৃত্তি করি,
  যাতে প্রত্যেক প্রযোজ্য বিধি শেষ পর্যন্ত বিবেচিত হয়।
  এভাবে নির্মিত !!{tableau}-র প্রত্যেক খোলা শাখা সম্পূর্ণ।
  % BN-SRC-381: উৎসে ‘বদ্ধ’ বলা; এখানে সম্পূর্ণ খোলা শাখার দাবি।
% BN-SRC-734 TARGET-CORRECTION-BEGIN
  \emph{উৎস-সংশোধন:} বদ্ধ শাখায় কাজ থামে; সেই শাখার সব সূত্র বিশ্লেষণের দাবি নেই। খোলা শাখার সম্পূর্ণতা সীমায় নেওয়া নির্মাণের ধর্ম, প্রতিটি সসীম পর্যায়ের ধর্ম নয়। প্রত্যেক পর্যায়ের শুরুতে থাকা সসীম সূত্র ও সূত্র--ব্যবহৃত-উত্তরসূরি জোড়ার সব অপূর্ণ কাজ একটি স্থির ক্রমে করি। নতুন সাক্ষী লাগে এমন প্রতিটি চিহ্নিত সূত্রের জন্য শাখায় একবার সাক্ষী নিই। নতুন সূত্র ও নাম পরের পর্যায়ের তালিকায় যায়। ফলে কোনো সূত্র বা প্রাসঙ্গিক উত্তরসূরি স্থায়ীভাবে বাদ পড়ে না।
  পরের উপপাদ্যের সম্ভবত অসীম $\Gamma$-র জন্যও সূত্রগুলি তালিকাবদ্ধ করা যায়: ভাষার সূত্রসমষ্টি গণনাযোগ্য। প্রতিটি পর্যায়ে পরের অনুমানটি সাধারণ মূল নাম $1$-এ প্রত্যেক খোলা শাখায় যোগ করি, তারপর ওই সসীম পর্যায়ের কাজ করি। কোনো পর্যায়ে সব শাখা বন্ধ হলে ইতিমধ্যে যোগ করা সসীমসংখ্যক অনুমান থেকেই একটি সসীম বদ্ধ ট্যাবলো পাওয়া যায়। এমন বদ্ধ ট্যাবলো না থাকলে প্রত্যেক পর্যায়ে খোলা শাখা থাকে। পর্যায়গুলি নিয়ে গড়া সসীম-শাখনবিশিষ্ট বৃক্ষে কনিগের সহায়ক উপপাদ্য প্রয়োগে একটি অবিরত খোলা পথ পাওয়া যায়; নির্মাণ শেষ হয়ে খোলা শাখা থেকে গেলে সেটিই নিই। নিরপেক্ষ ক্রমের জন্য ওই শাখায় সব অনুমান ও প্রতিটি প্রযোজ্য বিধির প্রয়োজনীয় উপসংহার থাকে। অতএব সেটি খোলা ও সম্পূর্ণ। অসীম নির্মাণটি মডেল পাওয়ার প্রমাণ; কোনো অসীম বদ্ধ বৃক্ষকে সসীম প্রমাণ বলা হচ্ছে না।
% BN-SRC-734 TARGET-CORRECTION-END
\end{proof}

\begin{thm}[সম্পূর্ণতা]
  \ollabel{thm:tableau-completeness}
  $\Gamma$-র কোনো বন্ধ !!{tableau} না থাকলে
  $\Gamma$ পরিতৃপ্তিযোগ্য।
\end{thm}

\begin{proof}
% BN-SRC-735 TARGET-CORRECTION-BEGIN
শূন্য অনুমানসমষ্টি হলে একটি জগৎবিশিষ্ট যেকোনো মডেলই তাকে পরিতৃপ্ত করে। তাই এখন $\Gamma$ অশূন্য ধরি; নির্বাচিত শাখায় একটি অনুমান আছে, ফলে $P(\Delta)$-ও অশূন্য এবং নিচের নির্মাণ সত্যিই একটি মডেল।
% BN-SRC-735 TARGET-CORRECTION-END
প্রস্তাবনা অনুযায়ী, $\Gamma$-র এমন একটি !!a{tableau} আছে যার
প্রত্যেক খোলা শাখা সম্পূর্ণ। এর কোনো বন্ধ !!a{tableau}
নেই বলে এতে অন্তত একটি শাখা খোলা ও সম্পূর্ণ।
সেই শাখার পূর্বসূচিযুক্ত !!{formula}s-এর সেট $\Delta$,
আর তাতে উপস্থিত পূর্বসূচিগুলির সেট $P(\Delta)$ ধরি।

$\mModel{M(\Delta)} = \tuple{P(\Delta), R, V}$ মডেলটি
নির্ধারণ করি। এর জগৎগুলি $\Delta$-তে উপস্থিত পূর্বসূচি;
অভিগম্যতা-সম্পর্কটি:
\[
R\sigma\sigma' \quad \text{যদি এবং কেবল যদি} \quad
\sigma'=\sigma.n \quad \text{কোনো~$n$-এর জন্য}
\]
এবং
\[
V(p) = \Setabs{\sigma}{\sFmla{\True}{p}[\sigma] \in \Delta}.
\]
$!A$-এর গঠনের উপর আরোহের সাহায্যে দেখাব,
$\sFmla{\True}{!A}[\sigma] \in \Delta$ হলে
$\mSat{M(\Delta)}{!A}[\sigma]$, আর
$\sFmla{\False}{!A}[\sigma] \in \Delta$ হলে
$\mSat/{M(\Delta)}{!A}[\sigma]$।

\begin{enumerate}
  \item \indcase{!A}{p}{$\sFmla{\True}{\indfrm}[\sigma] \in \Delta$
    হলে $V$-এর সংজ্ঞা অনুসারে $\sigma \in V(p)$; তাই
    $\mSat{M(\Delta)}{\indfrm}[\sigma]$.

    $\sFmla{\False}{\indfrm}[\sigma] \in \Delta$ হলে
    $\sFmla{\True}{\indfrm}[\sigma] \notin \Delta$,
    নইলে শাখাটি বন্ধ হতো। অতএব $\sigma \notin V(p)$ এবং
    $\mSat/{M(\Delta)}{\indfrm}[\sigma]$.}
  \item \indcase{!A}{\lnot !B}{\iftag{probNot}{অনুশীলন।}{
      $\sFmla{\True}{\indfrm}[\sigma] \in \Delta$ হলে
      শাখা সম্পূর্ণ বলে $\sFmla{\False}{!B}[\sigma] \in \Delta$।
      আরোহ অনুমানে $\mSat/{M(\Delta)}{!B}[\sigma]$; তাই
      $\mSat{M(\Delta)}{\indfrm}[\sigma]$.

      $\sFmla{\False}{\indfrm}[\sigma] \in \Delta$ হলে
      শাখা সম্পূর্ণ বলে $\sFmla{\True}{!B}[\sigma] \in \Delta$।
      আরোহ অনুমানে $\mSat{M(\Delta)}{!B}[\sigma]$; তাই
      $\mSat/{M(\Delta)}{\indfrm}[\sigma]$.}}
  \item \indcase{!A}{!B \land !C}{\iftag{probAnd}{অনুশীলন।}{
      $\sFmla{\True}{\indfrm}[\sigma] \in \Delta$ হলে
      শাখা সম্পূর্ণ বলে $\sFmla{\True}{!B}[\sigma] \in \Delta$
      এবং $\sFmla{\True}{!C}[\sigma] \in \Delta$।
      আরোহ অনুমানে
      $\mSat{M(\Delta)}{!B}[\sigma]$ এবং
      $\mSat{M(\Delta)}{!C}[\sigma]$। তাই
      $\mSat{M(\Delta)}{\indfrm}[\sigma]$.

      $\sFmla{\False}{\indfrm}[\sigma] \in \Delta$ হলে
      শাখা সম্পূর্ণ বলে
      $\sFmla{\False}{!B}[\sigma] \in \Delta$ অথবা
      $\sFmla{\False}{!C}[\sigma] \in \Delta$। আরোহ অনুমানে
      $\mSat/{M(\Delta)}{!B}[\sigma]$ অথবা
      % BN-SRC-382: দ্বিতীয় মিথ্যা সংযুক্ত উপসূত্রটি !C।
      $\mSat/{M(\Delta)}{!C}[\sigma]$। তাই
      $\mSat/{M(\Delta)}{\indfrm}[\sigma]$.}}
  \item \indcase{!A}{!B \lor !C}{\iftag{probOr}{অনুশীলন।}{
      $\sFmla{\True}{\indfrm}[\sigma] \in \Delta$ হলে
      শাখা সম্পূর্ণ বলে $\sFmla{\True}{!B}[\sigma] \in \Delta$
      অথবা $\sFmla{\True}{!C}[\sigma] \in \Delta$।
      আরোহ অনুমানে $\mSat{M(\Delta)}{!B}[\sigma]$
      অথবা $\mSat{M(\Delta)}{!C}[\sigma]$। তাই
      $\mSat{M(\Delta)}{\indfrm}[\sigma]$.

      $\sFmla{\False}{\indfrm}[\sigma] \in \Delta$ হলে
      শাখা সম্পূর্ণ বলে $\sFmla{\False}{!B}[\sigma] \in \Delta$
      এবং $\sFmla{\False}{!C}[\sigma] \in \Delta$।
      আরোহ অনুমানে $\mSat/{M(\Delta)}{!B}[\sigma]$
      এবং
      % BN-SRC-383: দ্বিতীয় মিথ্যা বিকল্প উপসূত্রটি !C।
      $\mSat/{M(\Delta)}{!C}[\sigma]$। তাই
      $\mSat/{M(\Delta)}{\indfrm}[\sigma]$.}}
  \item \indcase{!A}{!B \lif !C}{\iftag{probIf}{অনুশীলন।}{
      $\sFmla{\True}{\indfrm}[\sigma] \in \Delta$ হলে
      শাখা সম্পূর্ণ বলে $\sFmla{\False}{!B}[\sigma] \in \Delta$
      অথবা $\sFmla{\True}{!C}[\sigma] \in \Delta$।
      আরোহ অনুমানে $\mSat/{M(\Delta)}{!B}[\sigma]$
      অথবা $\mSat{M(\Delta)}{!C}[\sigma]$। তাই
      $\mSat{M(\Delta)}{\indfrm}[\sigma]$.

      $\sFmla{\False}{\indfrm}[\sigma] \in \Delta$ হলে
      শাখা সম্পূর্ণ বলে $\sFmla{\True}{!B}[\sigma] \in \Delta$
      এবং $\sFmla{\False}{!C}[\sigma] \in \Delta$।
      আরোহ অনুমানে $\mSat{M(\Delta)}{!B}[\sigma]$
      এবং
      % BN-SRC-384: মিথ্যা শর্তবাক্যের উত্তরপক্ষটি !C।
      $\mSat/{M(\Delta)}{!C}[\sigma]$। তাই
      $\mSat/{M(\Delta)}{\indfrm}[\sigma]$.}}
  \item \indcase{!A}{\Box !B}{\iftag{probBox}{অনুশীলন।}{
      $\sFmla{\True}{\indfrm}[\sigma] \in \Delta$ হলে
      শাখা সম্পূর্ণ বলে শাখায় ব্যবহৃত প্রত্যেক $\sigma.n$-এর
      জন্য $\sFmla{\True}{!B}[\sigma.n] \in \Delta$।
      অর্থাৎ, $R\sigma\sigma'$ মেনে চলা প্রত্যেক
      $\sigma' \in P(\Delta)$-র জন্যও তাই।
      আরোহ অনুমানে $R\sigma\sigma'$ মেনে চলা প্রত্যেক $\sigma'$-তে
      $\mSat{M(\Delta)}{!B}[\sigma']$। অতএব
      $\mSat{M(\Delta)}{\indfrm}[\sigma]$.

      $\sFmla{\False}{\indfrm}[\sigma] \in \Delta$ হলে
      শাখা সম্পূর্ণ বলে কোনো $\sigma.n$-এর জন্য
      $\sFmla{\False}{!B}[\sigma.n] \in \Delta$।
      আরোহ অনুমানে $\mSat/{M(\Delta)}{!B}[\sigma.n]$।
      যেহেতু $R\sigma(\sigma.n)$, এমন একটি
      $\sigma'$ আছে যার জন্য
      $\mSat/{M(\Delta)}{!B}[\sigma']$। তাই
      $\mSat/{M(\Delta)}{\indfrm}[\sigma]$.}}
  \item \indcase{!A}{\Diamond !B}{\iftag{probDiamond}{অনুশীলন।}{
      $\sFmla{\True}{\indfrm}[\sigma] \in \Delta$ হলে
      শাখা সম্পূর্ণ বলে কোনো $\sigma.n$-এর জন্য
      $\sFmla{\True}{!B}[\sigma.n] \in \Delta$।
      আরোহ অনুমানে $\mSat{M(\Delta)}{!B}[\sigma.n]$।
      যেহেতু $R\sigma(\sigma.n)$, এমন একটি
      $\sigma'$ আছে যার জন্য
      $\mSat{M(\Delta)}{!B}[\sigma']$। তাই
      $\mSat{M(\Delta)}{\indfrm}[\sigma]$.

      $\sFmla{\False}{\indfrm}[\sigma] \in \Delta$ হলে
      শাখা সম্পূর্ণ বলে শাখায় ব্যবহৃত প্রত্যেক $\sigma.n$-এর
      জন্য $\sFmla{\False}{!B}[\sigma.n] \in \Delta$।
      অর্থাৎ, $R\sigma\sigma'$ মেনে চলা প্রত্যেক
      $\sigma' \in P(\Delta)$-র জন্যও তাই।
      আরোহ অনুমানে $R\sigma\sigma'$ মেনে চলা প্রত্যেক $\sigma'$-তে
      $\mSat/{M(\Delta)}{!B}[\sigma']$। অতএব
      $\mSat/{M(\Delta)}{\indfrm}[\sigma]$.}}
\end{enumerate}
যেহেতু $\Gamma \subseteq \Delta$,
$\mSat{M(\Delta)}{\Gamma}$।
\end{proof}

\begin{prob}
\olref[nml][tab][cpl]{thm:tableau-completeness}-এর
প্রমাণ সম্পূর্ণ করুন।
\end{prob}

\begin{cor}
\ollabel{cor:entailment-completeness}
$\Gamma \Entails !A$ হলে $\Gamma \Proves !A$।
\end{cor}

\begin{cor}
\ollabel{cor:weak-completeness}
$!A$ সব মডেলে সত্য হলে $\Proves !A$।
\end{cor}

\end{document}
